\documentclass{article}

\usepackage{fancyhdr}
\usepackage{extramarks}
\usepackage{amsmath}
\usepackage{amsthm}
\usepackage{amsfonts}
\usepackage{tikz}
\usepackage[plain]{algorithm}
\usepackage{algpseudocode}
\usepackage{fullpage, url, hyperref}


\theoremstyle{plain}
\newtheorem{theorem}{Theorem}
\newtheorem{lemma}{Lemma}
\newtheorem{definition}{Definition}

\newcommand{\I}{\mathcal{I}}
\newcommand{\F}{\mathbb{F}}
\newcommand{\E}{\mathbb{E}}
\newcommand{\R}{\mathbb{R}}
\renewcommand{\P}{\mathbb{P}}
\usetikzlibrary{automata,positioning}

%
% Basic Document Settings
%

\topmargin=-0.45in
\evensidemargin=-.50in
\oddsidemargin=-.5in
\textwidth=7in
\textheight=9.5in
\headsep=0.25in

\linespread{1.1}

\pagestyle{fancy}
\lhead{\hmwkAuthorName}
\chead{\hmwkClass\ (\hmwkClassInstructor\ \hmwkClassTime): \hmwkTitle}
\rhead{\firstxmark}
\lfoot{\lastxmark}
\cfoot{\thepage}

\newcommand{\w}[2]{w_{#1}^{(#2)}}
\newcommand{\loss}[2]{\vec{\ell}_{#2}^{(#1)}}
\newcommand{\losst}[1]{\vec{\ell}^{(#1)}}
  \newcommand{\pt}[1]{\vec{p}^{(#1)}}
\renewcommand\headrulewidth{0.4pt}
\renewcommand\footrulewidth{0.4pt}

\setlength\parindent{0pt}


% 
% Create Problem Sections
%


\newcommand{\enterProblemHeader}[1]{
    \nobreak\extramarks{}{}\nobreak{}
    \nobreak\extramarks{}\nobreak{}
}

\newcommand{\exitProblemHeader}[1]{
    \nobreak\extramarks{}\nobreak{}
    \stepcounter{#1}
    \nobreak\extramarks{}\nobreak{}
  }

  
\newcommand{\enterExerciseHeader}[1]{
    \nobreak\extramarks{}{}\nobreak{}
    \nobreak\extramarks{}\nobreak{}
}

\newcommand{\exitExerciseHeader}[1]{
    \nobreak\extramarks{}\nobreak{}
    \stepcounter{#1}
    \nobreak\extramarks{}\nobreak{}
}

\setcounter{secnumdepth}{0}
\newcounter{partCounter}
\newcounter{homeworkProblemCounter}
\setcounter{homeworkProblemCounter}{1}

\newcounter{homeworkExerciseCounter}
\setcounter{homeworkExerciseCounter}{1}
\nobreak\extramarks{}\nobreak{}
\nobreak\extramarks{}\nobreak{}

%
% Homework Problem Environment
%
% This environment takes an optional argument. When given, it will adjust the
% problem counter. This is useful for when the problems given for your
% assignment aren't sequential. See the last 3 problems of this template for an
% example.
%
\newenvironment{homeworkProblem}[1][-1]{
%    \ifnum#1>0
   %     \setcounter{homeworkProblemCounter}{#1}
   % \fi
    \section{Problem \arabic{homeworkProblemCounter}}
    \setcounter{partCounter}{1}
    \enterProblemHeader{homeworkProblemCounter}
}{
    \exitProblemHeader{homeworkProblemCounter}
}

  \newenvironment{homeworkExercise}[1][-1]{
    \ifnum#1>0
        \setcounter{homeworkExerciseCounter}{#1}
    \fi
    \section{Exercise \arabic{homeworkExerciseCounter}}
    \setcounter{partCounter}{1}
    \enterExerciseHeader{homeworkExerciseCounter}
}{
    \exitExerciseHeader{homeworkExerciseCounter}
}

%
% Homework Details
%   - Title
%   - Due date
%   - Class
%   - Section/Time
%   - Instructor
%   - Author
%

\newcommand{\hmwkTitle}{Homework\ \#2}
\newcommand{\hmwkAssignedDate}{February 12, 2017}
\newcommand{\hmwkDueDate}{February 26, 2017, before class}
\newcommand{\hmwkClass}{Design and Analysis of Algorithms, CS 6550}
\newcommand{\hmwkClassTime}{}%Section A}
\newcommand{\hmwkClassInstructor}{Professor Jamie Morgenstern}
\newcommand{\hmwkAuthorName}{}%\textbf{Josh Davis} \and \textbf{Davis Josh}}

%
% Title Page
%

\title{    \vspace{2in}
    \textmd{\textbf{\hmwkClass:\ \hmwkTitle}}\\
    \normalsize\vspace{0.1in}\small{Due\ on\ \hmwkDueDate\ at 3:10pm}\\
    \vspace{0.1in}\large{\textit{\hmwkClassInstructor\ \hmwkClassTime}}
    \vspace{3in}
}

\author{\hmwkAuthorName}
\date{}

\renewcommand{\part}[1]{\textbf{\large Part \Alph{partCounter}}\stepcounter{partCounter}\\}

%
% Various Helper Commands
%

% Useful for algorithms
\newcommand{\alg}[1]{\textsc{\bfseries \footnotesize #1}}

% For derivatives
\newcommand{\deriv}[1]{\frac{\mathrm{d}}{\mathrm{d}x} (#1)}

% For partial derivatives
\newcommand{\pderiv}[2]{\frac{\partial}{\partial #1} (#2)}

% Integral dx
\newcommand{\dx}{\mathrm{d}x}

% Alias for the Solution section header
\newcommand{\solution}{\textbf{\large Solution}}

% Probability commands: Expectation, Variance, Covariance, Bias
\newcommand{\Var}{\mathrm{Var}}
\newcommand{\Cov}{\mathrm{Cov}}
\newcommand{\Bias}{\mathrm{Bias}}
\usepackage{enumitem}
\newcommand\litem[1]{\item{\bfseries#1.\space}}


\begin{document}

% \maketitle

\pagebreak

{\bf Homework Out: \hmwkAssignedDate}\\
{\bf Due Date:  \hmwkDueDate}\\

The HW contains some exercises (fairly simple problems to check you
are on board with the concepts; don’t submit your solutions), and
problems (for which you should submit your solutions, and which will
be graded). Some problems have sub-parts that are exercises.  For this
problem set, it’s OK to work with others. (Groups of 2, maybe 3 max.)
That being said, please think about the problems yourself before
talking to others. Please cite all sources you use, and people you
work with. The expectation is that you try and solve these problems
yourself, rather than looking online explicitly for
answers. Submissions due at beginning of class on the due date.
Please check the Piazza for details on submitting your \emph{LaTeXed}
solutions.

\paragraph{Exercises}
\begin{enumerate}

\litem{Flows, Kings, Halls, and Colors} Recall that
  Konig's theorem says for a bipartite graph $G$, the size of the
  maximum cardinality of a matching in $G$ is equal to the size of the
  minimum vertex cover.
  \begin{enumerate}
  \item Use the max-flow/min-cut theorem to prove Konig's
    theorem. (Recall that the max-flow/min-cut theorem says that in
    any flow network, the maximum $s-t$ flow equals the minimum $s-t$
    cut. Moreover if the arc capacities are integers, then there is a
    max flow which is integral.)
\item Use Konig’s theorem to prove Hall's theorem:\\
In a bipartite graph $G = (L, R, E)$, for any set $S \subseteq L$, let $N(S) = \{r \in R | \exists \ell \in
S,(\ell, r) \in E\}$ be the neighbors of $S$. Then $G$ has a matching of size $|L|$ if and
only if $|N(S)| \geq |S|$ for all $S \subseteq L$
\item Use Konig's or Hall’s theorems to show that in a bipartite graph where every vertex has
  the same degree $d$ (also called a $d$-regular bipartite graph), there is a perfect matching.
  \item Use the above part to show that the edges of a $d$-regular bipartite graph can be colored
using $d$ colors so that all edges incident on each vertex have different colors.
\end{enumerate}

\litem{Some Matching Reductions} Suppose you have an algorithm that solves max-weight
perfect matchings (MaxWPM) for all non-negative weight functions. Give reductions that
allow you to solve (a) min-weight perfect matchings (MinWPM) and (b) min-weight max cardinality
matchings (MinWMaxM)—i.e., among all matchings of size equal to $MM(G)$, find
the one with least weight.
If the MinWPM and MinWMaxM instances are bipartite, ensure that the reductions give you
MaxWPM instances that are bipartite too.

\litem{ Min-weight Perfect Matchings} Suppose $G = (L, R, E)$ is a undirected bipartite graph
with (possibly negative) integer edge weights $w_e$, suppose $M$ is some perfect matching. Let
$H_M$ be the digraph obtained by directing edges in $E \setminus M$ from left-to-right and putting weights
$w_e$ on them, and directing all edges in $M$ from right-to-left and putting weights $-w_e$ on them.
\begin{enumerate}
  \item  Show that $H_M$ has a negative-weight cycle if and only if $M$ is not a min-weight perfect
    matching.
    \item  Consider the algorithm: Start with any perfect matching $M$. While there exists a
negative-weight cycle $C$ in $H_M$, set $M \gets  M \Delta C$. Show that you will eventually get a
MwPM.

If $T(n)$ is the time to find a negative-weight cycle in an $n$-node graph, and $W := max_{e\in E} |w_e|$,
a naive bound on the runtime of the above algorithm would be $O(nW) \cdot T(n)$, plus the time
to find the first perfect matching.
\item  For matching $M$, denote $w(M) := \sum_{e \in M}w_e$. Show that if $w(M) > w(M^∗)$ then there
exists a cycle in $H_M$ with weight-ratio1 at most $\frac{w(M^∗)-w(M)}{n}$, where $M^*$ is  a min-weight
perfect matching.
\item  Suppose we change the algorithm in (b) to use the minimum weight-ratio cycle $C$ at
each step. Bound the runtime of this algorithm by $O(n \log(nW)) \cdot T'(n)$, plus the time
to find the first perfect matching. Here $T'(n)$ is the time to find a minimum weight-ratio
cycle.
\end{enumerate}

\litem{A Degree of Smoothness.} In the Max-Cut problem, we are given a graph $G = (V, E)$
with edge weights $w_e$, and want to find a subset $S \subseteq V$ such that the weight of the edges
crossing the cut (denoted as $\partial S := \{(u, v) \in E : u \in S, v \in V \setminus S\}$) is maximized. In the
smoothed analysis setting, the edge weights are random variables $W_e$ supported on $[0, 1]$,
whose density function is bounded by $1/\sigma$, as in lecture. Suppose the maximum degree in $G$
is $\Delta$.

One algorithm for max-cut is the simple local-search algorithm: start with any subset $S$.
Now if there is a vertex $v$ which can be moved from $S$ to $V \setminus S$ (or vice versa) to improve the
weight $w(\delta S)$, make the change. Show that the expected number of improving steps in such
an algorithm is bounded by $O(n^{3}2^\Delta/\sigma)$. (Hint: recall the argument we used for the $2OPT$
local search heuristic for TSP.)

\litem{Blocking Experts (due to Avrim Blum)} Here is a variation on the deterministic WeightedMajority
algorithm, designed to make it more adaptive.
\begin{itemize}
\item Each expert begins with weight 1 (as before).
\item We predict the result of a weighted-majority vote of the experts (as before).
\item If an expert makes a mistake, we penalize it by dividing its weight by 2, but only if its
weight was at least $\frac{1}{4}$ of the average weight of experts.
  \end{itemize}
Prove that in any contiguous block of trials (e.g., the 51st day through the 77th day), the
number of mistakes made by the algorithm is at most $O(m + \log N)$, where $m$ is the number
of mistakes made by the best expert \emph{in that block}, and $N$ is the total number of experts.

\litem{Expertise in Experts.} Here are some problems to help you internalize the experts proof.
\begin{enumerate}
  \item  Show that if we change the basic weighted majority algorithm to reduce the weight by
    $(1 − \epsilon)$ instead of $1/2$, we get
    \[\textrm{our total loss} \leq   2(1 + \epsilon) \cdot (\textrm{total loss of expert }i) + O(\frac{\log N}{\epsilon}).\]

    \item Show that if the loss vectors are in $[-\infty, 1]^N$ , then our analysis of Hedge gives us that
for any expert $i \in [N]$,
\[\sum _ { t } ( \vec { p } ^ { ( t ) } ,\vec { \ell } ^ { ( t ) } ) \leq \sum _ { t } \ell _ { i } ^ { ( t ) } + \frac { \ln N } { \epsilon } + \epsilon \sum _ { t } ( \vec { p } ^ { ( t ) } ,( \vec { \ell } ^ { ( t } ) ^ { 2} )\]
%
Here $\langle u, v \rangle$ is the usual inner product between two vectors $u$, $v$, and $(\vec{u})^2$ denotes a vector
whose $j$th entry is $u^2_j$.

\item Use the previous part to show that:
  \begin{enumerate}
  \item If losses are in $[0, 1]^N$ and $\epsilon \leq \frac{1}{2}$, then infer that our loss is
    \[\sum _ { t } ( \vec { p } ^ { ( t ) } ,\vec { \ell } ^ { ( t ) } \rangle \leq \frac { 1} { 1- \epsilon } ( \sum _ { t } \ell _ { i } ^ { ( t ) } + \frac { \ln N } { \epsilon } ) \leq ( 1+ 2\epsilon ) \sum _ { t } \ell _ { i } ^ { ( t ) } + \frac { 2\ln N } { \epsilon }.\]
    This gives a nice multiplicative-additive guarantee. (Hint: $x^2 \leq x$ for $x \in [0,1]$).
  \item Similarly, if we have only gains (i.e., if losses are all in  $[−1, 0]^N$ ), then we get
    \[text{ our gain } \geq ( 1- 2\epsilon ) ( \text{ gain of best expert } ) - \frac { 2\ln N } { \epsilon }.\]
  \end{enumerate}
\item Suppose the loss functions $\ell^{(t)} \in [- \rho, \rho]^N$, considver the vectors $p^{(t)}$ obtaine by running Hedge with ``learning rate'' $\epsilon^\prime = \frac{\epsilon}{2\rho}$ where we feed it the ``fake'' loss function $\tilde{\ell}^{(t)} - \vec\ell^{t}/\rho \in [-1,1]^N$. Show they satisfy
\[\sum _ { t } ( \vec { p } ^ { ( t ) } ,\vec { \ell } ^ { ( t ) } \rangle \leq \sum _ { t } \ell _ { i } ^ { ( t ) } + \frac { 2\rho ^ { 2} \ln N } { \epsilon } + \epsilon T / 2\]
 % 
    Hence, show that the average regret is at most $\epsilon$ if $T \geq \frac{2 \rho^2 \ln N}{\epsilon^2}$.
  \item Suppose we change the update step in the Hedge algorithm to be $\w{j}{t+1} \gets \w{j}{t} \cdot (1 - \epsilon \loss{j}{t})$.
    Alter the analysis to show that the regret after $T$ steps is still $O\left(\frac{\log N}{\epsilon} + \epsilon T\right)$.
  
\end{enumerate}
\end{enumerate}

\paragraph{Problems}

Answer problems \#1 and \#2, and any two out of the remaining three problems.

\renewcommand{\I}{\mathcal{I}}
\newcommand{\M}{\mathcal{M}}
\newcommand{\A}{\mathcal{A}}
\begin{enumerate}
  \litem{Matchings and Matroids} Let $\M = (U, \I)$ be a matroid. For
  $X \in \I$, define a bipartite graph
  $H_\M(X) = (X, U\setminus X, E)$, where $(x,y) \in E$ iff
  $x \in X, y \in U\setminus X$ and $(X + y - x)\in \I$. (that is, we
  can exchange $x$ for $y$ and get another independent set of the same
  size).
\begin{enumerate}
\item Let $X, Y \in \I$ be two independent sets with $|X| = |Y|$ Show
  that $H_\M(X)$ contains a perfect matching on $X \Delta Y$-- it
  contains a matchin $M \subseteq E$ which matches exactly the
  elements of $X \setminus Y$ to those of $Y \setminus X$ (Hint:
  Hall's theorem from the exercises).
\item (A partial converse.) For $X \in \I$, suppose there exists a set
  $Y\subseteq U$ such that $h-|m(X)$ contains a \emph{unique} perfect
  matching on $X \Delta Y$. Show that $Y \in \I$. HInt: one proof is
  via induction on $|Y\setminus X|$. You may use part (a).
  \end{enumerate}
%
\litem{Boost your score} Suppose you have a set of points $S \subseteq \R^d$, and for each point $x\in S$
you are also given a weight $w(x) \geq 1$ and a label $\ell(x) \in \{0, 1\}$. A function $f : \R^d \to \{0, 1\}$ is
called a hypothesis. If $W := \sum_{x \in S}w(x)$, the error of the hypothesis on the set $S$ is
\[\mathrm{err} _ { S ,l ,w } ( f ) : = \frac { 1} { W } \sum _ { x \in S } w ( \mathbf { x } ) \cdot 1_ { ( f ( x ) \neq \ell ( x ) ) },\]
i.e., the fraction of weight on points in $S$ that $f$ labels incorrectly.

Suppose you know an algorithm $\A$ that is “weak learner” for class $G$: given any set $S \subseteq \R^d$
(along with its weights $w()$ and the labels $\ell()$), this weak learner will return a function $f \in G$,
such that $\textrm{err}_{S,\ell,w}(f) \leq \frac{1}{2} -\eta$ for some $\eta > 0$. If $\eta$ is small, this means the returned hypothesis
is only slightly better than a random coin toss.

Give an algorithm that, given the point set $S$, along with the
weights $w()$ of the points and their labeling $\ell()$, runs $\A$ for
$T = O(\frac{1}{\eta^2} \log \frac{1}{\epsilon})$, times and
outputs a function $\hat{f}$ of the form
\[\hat { f } ( x ) = \mathbf{ m a j o r i t y ~ } ( f _ { 1} ( x ) ,f _ { 2} ( x ) ,\ldots ,f _ { T } ( x ) )\]

for $f_i \in G$, such that
$\textrm{err}_{S, \ell, w}(\hat{f}) \leq \epsilon$. (Hint: Start off
with the simple case where the weights of the points in $S$ are
$w(x) = 1$ for all $x \in S$; it contains all the ideas. Also, recall
the analysis of the Hedge or Weighted Majority even if you don’t use
these algorithms themselves.)

\litem{A Market-Based Bipartite Max-Matching Algorithm} Give $G = (I, B, E)$, the left vertices are items, the right are buyers. Let $n = \max(|I|, |B|)$. For each edge $(i,b) \in E$, let $v_{ib} = 1$; if $(i,b)\not\in E$, then $v_{ib} = 0$. The intial prices are $p_i = 0$; define the utility of buyer $b$ for item $i$ under prices $\overline{p} = (p_1, p_2 \ldots, p_{|I|})$ to be
\[u _ { i b } ( \overline { p } ) = m a x \{ v _ { i b } - p _ { i } ,0\}.\]
%
For some parameter $\delta \in (0,1)$, consider the following algorithm:
\begin{itemize}
\item{A1} Start with the empty matching $M = \emptyset$.
\item{A2} Pick any unmatched buyer $b$ such that its highest-utility
  item $i$ has $u_{ib}(\overline{p})\geq \delta$. Match $(i,b)$, which
  ma require dropping $(i,b')$ for some other $b'$ from the current
  matching.  After this operation, $i$ is assigned to $b$ rather than
  $b'$. Raise the price $p_i \gets p_i + \delta$.
\item{A3} If for each unmatched buyer,
  $\max_i u_{ib}(\overline{p}) < \delta$, let the current matching be
  denoted by $M_1$. Use the augmenting paths algorithm to augment
  $M_1$ to a maximum matching $M^*$.

  We now show that $M_1$ is a ``large'' matching, and so step $(A3)$ does not take ``much'' time.
\end{itemize}
\begin{enumerate}
\item Show that $|M_1| \geq |M^*| - n\delta$. Hint: you could try to use the following LP:
  \begin{align*}
    \left.\begin{array} { ll } { \max \sum _ { i b } v _ { i b } x _ { i b } } \\
            { \sum _ { b } x _ { i b } \leq 1} & { \forall i \in I } \\
            { \sum _ { i } x _ { i b } \leq 1} & { \forall b \in B } \end{array} \right.
  \end{align*}
  and show that $|M_1| \geq OPT_{LP} -  n \delta$ using the duals.
\item Show how to use the appropriate data structures to implement the algorithm so that all executions of step (A2) can be done in a tota of $O(\frac{1}{\delta}\cdot (m + n \ln n))$ time.
\item Choose the value of $\delta$ so thta the running time of the
  entire algorithm is as close to $O(m \sqrt{n})$ as possible. You may
  assume that each augmenting path can be found in $O(m)$ time and
  hence (A3) takes a total of $O(m) \times (|M^*| - |M_1|)$ time.

  Observe that this algorithm is different from the one in class:
  there the prices went up in lock-step, here we just pick a single
  item and unilaterally raise its price by a tiny constant. We also
  don’t get a perfect matching here at the end, we have to do some
  correction (Step (A3)).  BTW, what can you say about the performance
  of such an algorithm when you have general weights?
\end{enumerate}
\litem{Sherman-Morrison and Dynamic Algorithms}. Suppose graph $G_0$ has a perfect matching
to start off, and someone starts adding and removing edges. Each update step adds or
removes a single edge $e_i$ from $G_{i−1}$ to give $G_i$. At each time $i$, we want to answer ``does $G_i$
still have a perfect matching?” Of course, we could simply answer the question for each $G_i$
independently, but let us try to be smarter.

The Sherman-Morrison(-Woodbury) formula\footnote{See this
  \href{https://math.stackexchange.com/questions/252367/insightful-proofs-for-sherman-morrison-formula-and-matrix-determinant-lemma}{stackexchange
    discussion} for intuition, or the
  \href{https://en.wikipedia.org/wiki/Sherman\%E2\%80\%93Morrison_formula}{wikipedia article} for other pointers.}  says: for any non-singular
$n\times n$ matrix $A$, and any $n$-dimensional vectors $u, v$,
\begin{align*}
  ( A + u v ^ { T } ) ^ { - 1} = A ^ { - 1} - \frac { A ^ { - 1} u v ^ { T } A ^ { - 1} } { 1+ v ^ { T } A ^ { - 1} u .}
\end{align*}
(Assuming that $1 + v^T A^{-1} u \neq 0$). Moreover, the
\href{http://en.wikipedia.org/wiki/Matrix_determinant_lemma}{matrix
  determinant lemma} says
\begin{align*}
\mathrm{det} ( A + u v ^ { T } ) = \operatorname{det} ( A ) ( 1+ v ^ { T } A ^ { - 1} u )
\end{align*}

These tell us how the inverse and determinant of a matrix change under rank-one updates,
as long as these updates don’t make the matrices singular.
\begin{enumerate}
\item Use these to show that for $i \geq 1$, if
  $G_0, G_1\ldots, G_{i-1}$ all have perfect matchings, we can answer
  the question for $G_i$ correctly with probability
  $1 − O(i/poly(n))$, but in time $\tilde{O}(n^2)$. (Be sure to
  mention what field $\mathbb{F}_q$ you are working
  over.)\footnote{The naive way would be to use the Tutte-Lovasz
    approach for each timestep independently, which would require
    $O(n^\omega) \gg O(n^2)$ time per update: your goal is to do
    better} You may assume that the graphs $G_i$ are all bipartite.

  (As soon as $G_{i-1}$ has no perfect matching, your algorithm may be incorrect for all subsequent
$G_j , j \geq i$. Bonus: suggest ways to handle this shortcoming. Also bonus: suggest ways to
reduce remove the dependence on $i$ in the above probability bound.)

This general idea of using Sherman-Morrison for rank-one updates gives dynamic algorithms
for several other problems. Here is another example, where we need to count the number of
paths between vertices in DAGs.

\item  For a square matrix $M$ such that $M^k = 0$ for some positive integer $k$, show that
  \begin{align*}
( I - M ) ^ { - 1} = \sum _ { i = 0} ^ { k - 1} M ^ { i }.
  \end{align*}

  \item For a DAG $G$, let $A$ be the $n times n$ adjacency matrix: $A_{ij} = 1$ if and only if $(i, j)$ is an arc
in $G$. Show that the $(i, j)$
th entry of $(I − A)^{-1}$
is the number of paths between vertices
$i, j$ in $G$. Hence, infer that you can compute the number of paths from $i$ to $j$, for all $i, j$,
in time $O(n^\omega)$.
\item Show how to maintain these path counts under edge additions and deletions to the
graph, as long as it remains a DAG. (Make sure you argue that you don’t perform any
operations that cause matrices to become singular, etc.) The time for each update should
be $O(n^2)$. 
  \end{enumerate}

  \litem{Isolated? No Man is an Island} For this problem we assume that $G = (U \cup V, E)$ is a
bipartite graph with a perfect matching and a weight on the edges $w : E \to \R_{\geq 0}$. We say
that $w$ is “isolating” if the minimum-weight perfect matching in $G$ is unique.

You may use the following lemma:
\begin{lemma}
  If the edge weights are picked independently
  from a set of size $\geq 2|E| = 2m$, the resulting weight function is isolating with probability $\geq \frac{1}{2}$.
\end{lemma}
This requires $m \cdot \lceil \log_2{2m}\rceil$ bits of information. In
  this exercise, we construct an isolating weight with $O(\log^2 n)$
  bits of randomness.
  \begin{enumerate}
  \item (Do not submit) Given cycle $C = \langle u_1, v_1, u_2, v_2, . . . , u_k, v_k, u_1 \rangle$  in $G$, define its imbalance:
    \begin{align*}
i m b _ { w } ( C ) : =  w _ { u _ { 1} v _ { 1} } - w _ { v _ { 1} u _ { 2} } + w _ { u _ { 2} v _ { 2} } - w _ { v _ { 2} u _ { 3} } \cdots + w _ { u _ { k } v _ { k } } - w _ { v _ { k } u _ { 1} } |.
    \end{align*}
    Show that if $w$ is a weifht satisfying $imb_w(C) \neq 0$ for all cycles $C$, the $w$ is isolating.
  \item Let $h$ be a graph with $n$ nodes and no cycles of length
    $\leq r$ and let $r^\prime = 4 \cdot
    \lfloor r/2 \rfloor$. Then
    show that $H$ has at most $n^4$ cycles of length $\leq r^\prime$.
  \item Let $C$ be a cycle in $G$ such that $imb_w(G) \neq 0$. Let
    $E_1$ be the union of all minimum-weight perfect matchings in
    $G$. Show then that $G_1 = (U \cup V, E_1)$ does not contain a
    cycle in $C$.

    Hint: Consider the perfect matching polytope and let $x$ denote
    the point which is the average of all the min-weight perfect
    matchings in $G$. Show that if $imb_w(C) \neq 0$ then you can
    construct a new point that is cheaper than $x$.
  \item Consider a graph $G$ with $n$ nodes and some $s\geq 1$. You may assume that there is an oracle $O$
which uses $\log(ns)$ random bits to produce a non-negative weight $w$ such that for any
set of s cycles $\{C_1, C_2,\ldots ,C_s\}$,

\begin{align*}
\P[imb_w ( C _ { i } ) \neq 0\text{ for all }1 \leq i \leq s ] \geq 1- \frac { 1} { n }.
\end{align*}

Use the above oracle $O$ along with parts (b) and (c), to construct an isolating weight $w$
using $O(\log^2 n)$ bits of randomness with probability at least $1 -\frac{O(\log n)}{n}$.

    \end{enumerate}
   
\end{enumerate}



\end{document}
